\chapter{Subnetting and IPv6 Drills}\label{app:subnetting}

Subnetting appears on the exam mostly \emph{inside} other questions --- an ACL
wildcard, a summary route, whether two hosts share a subnet. It is rarely the
point of the question, which means every second you spend calculating is a second
taken from the question you were actually asked.

\begin{alertbox}[The standard to aim for]
Any \cmd{/24} to \cmd{/30} question, answered in \textbf{under fifteen seconds,
without writing anything down}. That is achievable by almost anyone with about
two hours of drilling, and it is worth more marks than any single chapter of this
book.

The drills in this appendix have their answers computed rather than typed. If
your answer differs from one here, work it through again --- the table is right.
\end{alertbox}

\section{The method}

Forget binary conversion. There is one number to find.

\begin{definitionbox}[The block size, or ``magic number'']
\textbf{256 minus the interesting octet of the mask.} Subnets begin at multiples
of that number, in that octet.
\end{definitionbox}

\begin{worked}[192.168.10.77/26]
\textbf{Step 1 --- the mask.} \cmd{/26} is 24 bits plus 2, so the fourth octet
has 2 bits: $128+64 = 192$. Mask \cmd{255.255.255.192}. The fourth octet is the
interesting one.

\textbf{Step 2 --- the block size.} $256 - 192 = \mathbf{64}$.

\textbf{Step 3 --- count in blocks} until you pass the address. 0, 64, 128, 192.
The address ends in 77, which falls between 64 and 128, so the network is
\textbf{192.168.10.64}.

\textbf{Step 4 --- the rest follows.} The next block starts at 128, so the
broadcast is one below it: \textbf{192.168.10.127}. Usable hosts are everything
between: \textbf{.65 to .126}, which is $64 - 2 = \mathbf{62}$.

Four steps, no binary, and with the two tables below memorised it is arithmetic
you can do while reading the question.
\end{worked}

\section{The two tables to memorise}

\begin{longtable}{@{}L{2.4cm}L{2.4cm}L{2.4cm}L{7.4cm}@{}}
\toprule
\rowcolor{primary!15}
\textbf{Mask octet} & \textbf{Block size} & \textbf{Bits} & \textbf{Subnets it makes in that octet} \\
\midrule
\endhead
\cmd{128} & 128 & 1 & 0, 128 \\
\rowcolor{lightbg}
\cmd{192} & 64 & 2 & 0, 64, 128, 192 \\
\cmd{224} & 32 & 3 & 0, 32, 64, 96, 128, 160, 192, 224 \\
\rowcolor{lightbg}
\cmd{240} & 16 & 4 & 0, 16, 32, 48, \ldots \\
\cmd{248} & 8 & 5 & 0, 8, 16, 24, \ldots \\
\rowcolor{lightbg}
\cmd{252} & 4 & 6 & 0, 4, 8, 12, \ldots \\
\cmd{254} & 2 & 7 & 0, 2, 4, 6, \ldots \\
\rowcolor{lightbg}
\cmd{255} & 1 & 8 & every value \\
\bottomrule
\end{longtable}

\begin{longtable}{@{}L{1.5cm}L{3.4cm}L{2.3cm}L{2.4cm}L{4.9cm}@{}}
\toprule
\rowcolor{primary!15}
\textbf{Prefix} & \textbf{Mask} & \textbf{Block} & \textbf{Usable} & \textbf{Typical use} \\
\midrule
\endhead
\cmd{/30} & \cmd{255.255.255.252} & 4 & 2 & A point-to-point link \\
\rowcolor{lightbg}
\cmd{/29} & \cmd{255.255.255.248} & 8 & 6 & A handful of servers \\
\cmd{/28} & \cmd{255.255.255.240} & 16 & 14 & A small office \\
\rowcolor{lightbg}
\cmd{/27} & \cmd{255.255.255.224} & 32 & 30 & A classroom \\
\cmd{/26} & \cmd{255.255.255.192} & 64 & 62 & A floor \\
\rowcolor{lightbg}
\cmd{/25} & \cmd{255.255.255.128} & 128 & 126 & Half a /24 \\
\cmd{/24} & \cmd{255.255.255.0} & 256 & 254 & The default assumption \\
\rowcolor{lightbg}
\cmd{/23} & \cmd{255.255.254.0} & 2 (3rd) & 510 & Two /24s \\
\cmd{/22} & \cmd{255.255.252.0} & 4 (3rd) & 1022 & A building \\
\rowcolor{lightbg}
\cmd{/21} & \cmd{255.255.248.0} & 8 (3rd) & 2046 & A campus block \\
\cmd{/20} & \cmd{255.255.240.0} & 16 (3rd) & 4094 & A site \\
\bottomrule
\end{longtable}

\begin{conceptbox}[Wildcard masks are the same table, inverted]
Subtract each octet of the mask from 255 (\chref{acl}). \cmd{/26} is
\cmd{255.255.255.192}, so the wildcard is \cmd{0.0.0.63} --- and 63 is one less
than the block size of 64, which is the quickest way to produce it.

\textbf{Check it every time:} network address plus wildcard should equal the
broadcast address. $192.168.10.64 + 0.0.0.63 = 192.168.10.127$. If it does not,
the wildcard is wrong.
\end{conceptbox}

\section{Drill 1 --- network, broadcast and range}

For each, give the mask in dotted decimal, the network address, the broadcast
address, the usable range and the host count. Answers in Section~\ref{sec:c-answers}.

\begin{longtable}{@{}L{0.8cm}L{4.2cm}L{3.1cm}L{3.1cm}L{3.3cm}@{}}
\toprule
\rowcolor{primary!15}
\textbf{} & \textbf{Given} & \textbf{Network} & \textbf{Broadcast} & \textbf{Usable hosts} \\
\midrule
\endhead
1 & \cmd{192.168.10.77/26} &  &  &  \\
\rowcolor{lightbg}
2 & \cmd{10.4.19.200/21} &  &  &  \\
3 & \cmd{172.16.35.14/20} &  &  &  \\
\rowcolor{lightbg}
4 & \cmd{192.168.1.130/25} &  &  &  \\
5 & \cmd{10.0.0.66/30} &  &  &  \\
\rowcolor{lightbg}
6 & \cmd{172.20.140.9/22} &  &  &  \\
7 & \cmd{192.168.200.35/27} &  &  &  \\
\rowcolor{lightbg}
8 & \cmd{10.55.12.3/19} &  &  &  \\
9 & \cmd{172.31.255.200/18} &  &  &  \\
\rowcolor{lightbg}
10 & \cmd{192.168.16.62/28} &  &  &  \\
11 & \cmd{10.128.64.31/27} &  &  &  \\
\rowcolor{lightbg}
12 & \cmd{172.16.0.129/29} &  &  &  \\
13 & \cmd{203.0.113.45/29} &  &  &  \\
\rowcolor{lightbg}
14 & \cmd{198.51.100.17/28} &  &  &  \\
15 & \cmd{192.0.2.200/26} &  &  &  \\
\rowcolor{lightbg}
16 & \cmd{10.10.10.10/23} &  &  &  \\
17 & \cmd{172.18.6.250/21} &  &  &  \\
\rowcolor{lightbg}
18 & \cmd{192.168.99.99/24} &  &  &  \\
\bottomrule
\end{longtable}

\section{Drill 2 --- sizing a subnet}

The other direction, and the one VLSM questions need: given a host requirement,
find the smallest prefix that fits. Work each out before reading the row.

\begin{longtable}{@{}L{3.4cm}L{3.0cm}L{3.4cm}L{4.6cm}@{}}
\toprule
\rowcolor{primary!15}
\textbf{Hosts needed} & \textbf{Prefix} & \textbf{Mask} & \textbf{Usable / wasted} \\
\midrule
\endhead
2 & \cmd{/30} & \cmd{255.255.255.252} & 2 / 0 \\
\rowcolor{lightbg}
6 & \cmd{/29} & \cmd{255.255.255.248} & 6 / 0 \\
12 & \cmd{/28} & \cmd{255.255.255.240} & 14 / 2 \\
\rowcolor{lightbg}
25 & \cmd{/27} & \cmd{255.255.255.224} & 30 / 5 \\
50 & \cmd{/26} & \cmd{255.255.255.192} & 62 / 12 \\
\rowcolor{lightbg}
100 & \cmd{/25} & \cmd{255.255.255.128} & 126 / 26 \\
200 & \cmd{/24} & \cmd{255.255.255.0} & 254 / 54 \\
\rowcolor{lightbg}
500 & \cmd{/23} & \cmd{255.255.254.0} & 510 / 10 \\
1000 & \cmd{/22} & \cmd{255.255.252.0} & 1022 / 22 \\
\bottomrule
\end{longtable}

\begin{alertbox}[The off-by-one that costs marks]
A requirement for 30 hosts needs a \cmd{/26}, not a \cmd{/27}. A \cmd{/27} has 32
addresses and \textbf{30 usable} --- so it fits exactly, with no room for the
gateway you forgot to count.

Count the gateway, and count any HSRP virtual address (\chref{fhrp}) as well: a
subnet with two routers running HSRP consumes three addresses before a single
user appears.
\end{alertbox}

\section{Drill 3 --- IPv6 compression}

Apply both rules: drop leading zeros in each group, and replace \emph{one} run of
consecutive all-zero groups with \cmd{::}.

\begin{longtable}{@{}L{0.8cm}L{7.6cm}L{6.9cm}@{}}
\toprule
\rowcolor{primary!15}
\textbf{} & \textbf{Given} & \textbf{Compressed} \\
\midrule
\endhead
1 & \cmd{2001:0db8:0000:0000:0000:ff00:0042:8329} &  \\
\rowcolor{lightbg}
2 & \cmd{2001:0db8:acad:0010:0000:0000:0000:0001} &  \\
3 & \cmd{fe80:0000:0000:0000:0204:61ff:fe9d:f156} &  \\
\rowcolor{lightbg}
4 & \cmd{2001:0db8:0000:0001:0000:0000:0000:0100} &  \\
5 & \cmd{0000:0000:0000:0000:0000:0000:0000:0001} &  \\
\rowcolor{lightbg}
6 & \cmd{2001:0db8:aaaa:0001:0000:0000:0000:0000} &  \\
\bottomrule
\end{longtable}

\begin{conceptbox}[The rule people get wrong]
\cmd{::} may appear \textbf{once}. \cmd{2001:db8::1::5} is not a valid address,
because a parser cannot tell how many zero groups each \cmd{::} stands for.

Where there are two runs of zeros, compress the \emph{longer} one and write the
shorter out in full. If they are equal in length, compress the first.
\end{conceptbox}

\section{Drill 4 --- EUI-64}

Build the full address from the prefix and the MAC. The procedure is: split the
MAC in half, insert \cmd{fffe} in the middle, and \textbf{flip the seventh bit of
the first octet}.

\begin{longtable}{@{}L{0.7cm}L{4.0cm}L{3.2cm}L{7.0cm}@{}}
\toprule
\rowcolor{primary!15}
\textbf{} & \textbf{Prefix} & \textbf{MAC} & \textbf{EUI-64 address} \\
\midrule
\endhead
1 & \cmd{2001:db8:acad:10::/64} & \cmd{00:1a:2b:3c:4d:5e} &  \\
\rowcolor{lightbg}
2 & \cmd{2001:db8:1:1::/64} & \cmd{aa:bb:cc:dd:ee:ff} &  \\
3 & \cmd{fe80::/64} & \cmd{0c:1d:2e:3f:40:51} &  \\
\rowcolor{lightbg}
4 & \cmd{2001:db8:99::/64} & \cmd{52:54:00:12:34:56} &  \\
\bottomrule
\end{longtable}

\begin{worked}[Why the seventh bit flips, and how to do it fast]
That bit is the universal/local flag. In a MAC address it is 0 for a
manufacturer-assigned address; EUI-64 inverts it, so a burned-in address becomes
1.

In practice you are flipping the value 2 in the second hex digit of the first
octet. \cmd{00} becomes \cmd{02}. \cmd{aa} becomes \cmd{a8}. \cmd{52} becomes
\cmd{50}. \cmd{0c} becomes \cmd{0e}.

If the second hex digit is even, add 2; if it is odd\ldots{} it will not be, for
a normal MAC. Learn the four cases above and you have covered nearly everything
you will be shown.
\end{worked}

\section{Answers}\label{sec:c-answers}

\subsection{Drill 1}

\begin{longtable}{@{}L{0.8cm}L{3.1cm}L{3.0cm}L{3.0cm}L{4.4cm}@{}}
\toprule
\rowcolor{primary!15}
\textbf{} & \textbf{Mask} & \textbf{Network} & \textbf{Broadcast} & \textbf{Usable range (count)} \\
\midrule
\endhead
1 & \cmd{255.255.255.192} & \cmd{192.168.10.64} & \cmd{192.168.10.127} & \cmd{192.168.10.65} \newline \cmd{to 192.168.10.126} (62) \\
\rowcolor{lightbg}
2 & \cmd{255.255.248.0} & \cmd{10.4.16.0} & \cmd{10.4.23.255} & \cmd{10.4.16.1} \newline \cmd{to 10.4.23.254} (2046) \\
3 & \cmd{255.255.240.0} & \cmd{172.16.32.0} & \cmd{172.16.47.255} & \cmd{172.16.32.1} \newline \cmd{to 172.16.47.254} (4094) \\
\rowcolor{lightbg}
4 & \cmd{255.255.255.128} & \cmd{192.168.1.128} & \cmd{192.168.1.255} & \cmd{192.168.1.129} \newline \cmd{to 192.168.1.254} (126) \\
5 & \cmd{255.255.255.252} & \cmd{10.0.0.64} & \cmd{10.0.0.67} & \cmd{10.0.0.65} \newline \cmd{to 10.0.0.66} (2) \\
\rowcolor{lightbg}
6 & \cmd{255.255.252.0} & \cmd{172.20.140.0} & \cmd{172.20.143.255} & \cmd{172.20.140.1} \newline \cmd{to 172.20.143.254} (1022) \\
7 & \cmd{255.255.255.224} & \cmd{192.168.200.32} & \cmd{192.168.200.63} & \cmd{192.168.200.33} \newline \cmd{to 192.168.200.62} (30) \\
\rowcolor{lightbg}
8 & \cmd{255.255.224.0} & \cmd{10.55.0.0} & \cmd{10.55.31.255} & \cmd{10.55.0.1} \newline \cmd{to 10.55.31.254} (8190) \\
9 & \cmd{255.255.192.0} & \cmd{172.31.192.0} & \cmd{172.31.255.255} & \cmd{172.31.192.1} \newline \cmd{to 172.31.255.254} (16382) \\
\rowcolor{lightbg}
10 & \cmd{255.255.255.240} & \cmd{192.168.16.48} & \cmd{192.168.16.63} & \cmd{192.168.16.49} \newline \cmd{to 192.168.16.62} (14) \\
11 & \cmd{255.255.255.224} & \cmd{10.128.64.0} & \cmd{10.128.64.31} & \cmd{10.128.64.1} \newline \cmd{to 10.128.64.30} (30) \\
\rowcolor{lightbg}
12 & \cmd{255.255.255.248} & \cmd{172.16.0.128} & \cmd{172.16.0.135} & \cmd{172.16.0.129} \newline \cmd{to 172.16.0.134} (6) \\
13 & \cmd{255.255.255.248} & \cmd{203.0.113.40} & \cmd{203.0.113.47} & \cmd{203.0.113.41} \newline \cmd{to 203.0.113.46} (6) \\
\rowcolor{lightbg}
14 & \cmd{255.255.255.240} & \cmd{198.51.100.16} & \cmd{198.51.100.31} & \cmd{198.51.100.17} \newline \cmd{to 198.51.100.30} (14) \\
15 & \cmd{255.255.255.192} & \cmd{192.0.2.192} & \cmd{192.0.2.255} & \cmd{192.0.2.193} \newline \cmd{to 192.0.2.254} (62) \\
\rowcolor{lightbg}
16 & \cmd{255.255.254.0} & \cmd{10.10.10.0} & \cmd{10.10.11.255} & \cmd{10.10.10.1} \newline \cmd{to 10.10.11.254} (510) \\
17 & \cmd{255.255.248.0} & \cmd{172.18.0.0} & \cmd{172.18.7.255} & \cmd{172.18.0.1} \newline \cmd{to 172.18.7.254} (2046) \\
\rowcolor{lightbg}
18 & \cmd{255.255.255.0} & \cmd{192.168.99.0} & \cmd{192.168.99.255} & \cmd{192.168.99.1} \newline \cmd{to 192.168.99.254} (254) \\
\bottomrule
\end{longtable}

\subsection{Drill 3}

\begin{longtable}{@{}L{0.8cm}L{7.6cm}L{6.9cm}@{}}
\toprule
\rowcolor{primary!15}
\textbf{} & \textbf{Given} & \textbf{Compressed} \\
\midrule
\endhead
1 & \cmd{2001:0db8:0000:0000:0000:ff00:0042:8329} & \cmd{2001:db8::ff00:42:8329} \\
\rowcolor{lightbg}
2 & \cmd{2001:0db8:acad:0010:0000:0000:0000:0001} & \cmd{2001:db8:acad:10::1} \\
3 & \cmd{fe80:0000:0000:0000:0204:61ff:fe9d:f156} & \cmd{fe80::204:61ff:fe9d:f156} \\
\rowcolor{lightbg}
4 & \cmd{2001:0db8:0000:0001:0000:0000:0000:0100} & \cmd{2001:db8:0:1::100} \\
5 & \cmd{0000:0000:0000:0000:0000:0000:0000:0001} & \cmd{::1} \\
\rowcolor{lightbg}
6 & \cmd{2001:0db8:aaaa:0001:0000:0000:0000:0000} & \cmd{2001:db8:aaaa:1::} \\
\bottomrule
\end{longtable}

\subsection{Drill 4}

\begin{longtable}{@{}L{0.7cm}L{4.0cm}L{3.2cm}L{7.0cm}@{}}
\toprule
\rowcolor{primary!15}
\textbf{} & \textbf{Prefix} & \textbf{MAC} & \textbf{EUI-64 address} \\
\midrule
\endhead
1 & \cmd{2001:db8:acad:10::/64} & \cmd{00:1a:2b:3c:4d:5e} & \cmd{2001:db8:acad:10:21a:2bff:fe3c:4d5e} \\
\rowcolor{lightbg}
2 & \cmd{2001:db8:1:1::/64} & \cmd{aa:bb:cc:dd:ee:ff} & \cmd{2001:db8:1:1:a8bb:ccff:fedd:eeff} \\
3 & \cmd{fe80::/64} & \cmd{0c:1d:2e:3f:40:51} & \cmd{fe80::e1d:2eff:fe3f:4051} \\
\rowcolor{lightbg}
4 & \cmd{2001:db8:99::/64} & \cmd{52:54:00:12:34:56} & \cmd{2001:db8:99:0:5054:ff:fe12:3456} \\
\bottomrule
\end{longtable}

\begin{alertbox}[Regenerating these]
Every answer above is computed by \cmd{tools/gen\_drills.py} using Python's
\cmd{ipaddress} module, not worked out by hand. To change or extend the drills,
edit the lists at the top of that script and re-run it --- do not edit the tables
here, because a hand-edited answer is exactly the error this arrangement exists
to prevent.
\end{alertbox}
